Los resultados experimentales obtenidos en este trabajo corroboran de manera contundente los modelos teóricos fundamentales que describen la propagación de la luz para fuentes puntuales y haces colimados. En el caso de la fuente puntual, representada por un LED, se verificó empíricamente la Ley Inversa al Cuadrado, lo que confirma su comportamiento como emisor isotrópico (frente de onda esférico) de radiación. La dependencia de la irradiancia con la distancia se ajustó a una relación potencial con un exponente de $-1.944 \pm 0.006$, un valor notablemente cercano al valor teórico de -2, lo que refuerza la validez del modelo y la precisión de las mediciones. En contraste, el haz colimado, generado por un láser, exhibió una irradiancia que se mantuvo prácticamente constante a lo largo de las distintas distancias de medición. Este comportamiento concuerda con la teoría de haces gaussianos y la mínima divergencia esperada en este tipo de fuentes. El ajuste de la irradiancia en función de la distancia reveló un exponente de $-0.0182 \pm 0.002$, lo que indica una muy ligera divergencia del haz real en comparación con lo que se le puede atribuir a una fuente de divergencia importante (como en este caso el led). Esta pequeña desviación del comportamiento ideal de un haz perfectamente colimado puede atribuirse a factores como las imperfecciones ópticas del sistema o los efectos de la difracción.

La comparación directa de la irradiancia normalizada para ambas fuentes ilustra de manera clara la diferencia fundamental en su comportamiento. Mientras que la fuente puntual muestra una rápida disminución de la irradiancia con la distancia, el haz colimado mantiene una irradiancia constante (al menos en nuestro intervalo de estudio que comprende desde los $(40\pm2)mm$ hasta los $(400\pm2)mm$), lo que destaca las ventajas de este último para aplicaciones que requieren la entrega de energía concentrada a distancias relativamente grandes. Los resultados obtenidos contribuyen al conocimiento fundamental de la óptica clásica y proporcionan herramientas valiosas para el desarrollo de nuevas tecnologías.

Como posible extensión de este trabajo, se podría considerar la investigación del comportamiento de las fuentes de luz a distancias mayores, así como el efecto de diferentes condiciones ambientales, como la temperatura o la humedad, en la propagación de la luz. Además, se podrían explorar otras características de las fuentes de luz, como la coherencia espacial y temporal, para obtener una caracterización más completa.